\documentclass[11pt,a4paper]{article}
\input{../template/utm-lab}

\newcommand{\labnumber}{11}
\newcommand{\labtitle}{The Race and the Semaphore}
\newcommand{\studentname}{Daniel Nicorici}
\newcommand{\studentgroup}{FAF-243}

\begin{document}

\utmtitlepage
\tableofcontents
\newpage

\plainsection{Purpose of the work}

Two threads that touch the same variable at the same time lose work. The purpose of this lab is
to see that happen, to fix it with a mutex and with a semaphore, and then to use the same
semaphore in three ways: as a lock, as a signal and as a counter.

The work was done on my Arch Linux laptop with gcc 16.2.1, in the folder
\code{\textasciitilde/os-lab11}. All programs were compiled with \code{-pthread} and without
\code{-O2}, because the optimizer would hide the race of Part 1.

\section{Part 1. See the race}

Two threads each add 1 to the same counter, one million times. The right answer is 2\,000\,000.

\begin{lstlisting}[style=ccode,caption={race.c}]
#include <stdio.h>
#include <pthread.h>

#define N 1000000

long counter = 0;

void *add(void *arg) {
    for (int i = 0; i < N; i++)
        counter++;
    return NULL;
}

int main(void) {
    pthread_t a, b;
    pthread_create(&a, NULL, add, NULL);
    pthread_create(&b, NULL, add, NULL);
    pthread_join(a, NULL);
    pthread_join(b, NULL);
    printf("counter = %ld (expected %d)\n", counter, 2 * N);
    return 0;
}
\end{lstlisting}

\begin{lstlisting}[caption={Three runs of race}]
$ gcc -Wall -pthread -o race race.c
$ ./race
counter = 1100695 (expected 2000000)
$ ./race
counter = 1022820 (expected 2000000)
$ ./race
counter = 1076745 (expected 2000000)
\end{lstlisting}

\textbf{Observations:}
\begin{enumerate}
  \item \textbf{The three steps of \code{counter++}.} The CPU loads the value of \code{counter}
  from memory into a register, adds 1 in the register, and stores the register back to memory.
  \item \textbf{How an addition is lost.} Both threads load the same value, for example 5. Each
  adds 1 in its own register and each stores 6. Two additions were done, but the counter grew
  only by one, so one of them is lost.
  \item \textbf{Why the result differs.} The moments at which the scheduler and the two CPU
  cores interleave the threads are different at every run, so a different number of additions
  collide each time. Almost half of the additions were lost in every run.
\end{enumerate}

\section{Part 2. Fix it: mutex, then semaphore}

\subsection{A mutex around counter++}

The line between lock and unlock is the critical section: only one thread may be inside.

\begin{lstlisting}[style=ccode,caption={mutex.c}]
#include <stdio.h>
#include <pthread.h>

#define N 1000000

long counter = 0;
pthread_mutex_t lock = PTHREAD_MUTEX_INITIALIZER;

void *add(void *arg) {
    for (int i = 0; i < N; i++) {
        pthread_mutex_lock(&lock);
        counter++;
        pthread_mutex_unlock(&lock);
    }
    return NULL;
}

int main(void) {
    pthread_t a, b;
    pthread_create(&a, NULL, add, NULL);
    pthread_create(&b, NULL, add, NULL);
    pthread_join(a, NULL);
    pthread_join(b, NULL);
    printf("counter = %ld (expected %d)\n", counter, 2 * N);
    return 0;
}
\end{lstlisting}

\subsection{The same lock, made from a semaphore}

A semaphore is a number that never goes below 0. \code{sem\_wait} takes one unit, and waits if
there is none; \code{sem\_post} gives one back. Started at 1, it behaves like a mutex.

\begin{lstlisting}[style=ccode,caption={sem.c}]
#include <stdio.h>
#include <pthread.h>
#include <semaphore.h>

#define N 1000000

long counter = 0;
sem_t s;

void *add(void *arg) {
    for (int i = 0; i < N; i++) {
        sem_wait(&s);
        counter++;
        sem_post(&s);
    }
    return NULL;
}

int main(void) {
    sem_init(&s, 0, 1);
    pthread_t a, b;
    pthread_create(&a, NULL, add, NULL);
    pthread_create(&b, NULL, add, NULL);
    pthread_join(a, NULL);
    pthread_join(b, NULL);
    printf("counter = %ld (expected %d)\n", counter, 2 * N);
    sem_destroy(&s);
    return 0;
}
\end{lstlisting}

\begin{lstlisting}[caption={mutex and sem, three runs each}]
$ gcc -Wall -pthread -o mutex mutex.c
$ ./mutex
counter = 2000000 (expected 2000000)
$ ./mutex
counter = 2000000 (expected 2000000)
$ ./mutex
counter = 2000000 (expected 2000000)
$ gcc -Wall -pthread -o sem sem.c
$ ./sem
counter = 2000000 (expected 2000000)
$ ./sem
counter = 2000000 (expected 2000000)
$ ./sem
counter = 2000000 (expected 2000000)
\end{lstlisting}

\subsection{The price of the lock}

\begin{lstlisting}[caption={Timing the three versions}]
$ time ./race
counter = 1074869 (expected 2000000)

real	0m0.004s
user	0m0.006s
sys	0m0.001s
$ time ./mutex
counter = 2000000 (expected 2000000)

real	0m0.032s
user	0m0.036s
sys	0m0.027s
$ time ./sem
counter = 2000000 (expected 2000000)

real	0m0.051s
user	0m0.069s
sys	0m0.033s
\end{lstlisting}

\begin{lstlisting}[caption={sem.c with sem\_init(\&s, 0, 2), three runs}]
$ gcc -Wall -pthread -o sem sem.c && ./sem && ./sem && ./sem
counter = 1978631 (expected 2000000)
counter = 1988948 (expected 2000000)
counter = 1982592 (expected 2000000)
\end{lstlisting}

\textbf{Observations:}
\begin{enumerate}
  \item \textbf{Always 2\,000\,000?} Yes. \code{mutex} and \code{sem} gave exactly 2\,000\,000
  in all three runs each.
  \item \textbf{A start value of 2.} The semaphore then lets two threads into the critical
  section at once, and the program has only two threads, so nobody ever waits and the race is
  back: 1\,978\,631, 1\,988\,948 and 1\,982\,592. Fewer additions are lost than in
  \code{race}, only because the extra calls make the collisions rarer. After the test the value
  was put back to 1.
  \item \textbf{Why the locked version is slower.} Every addition now also pays for one lock
  and one unlock operation, and a thread that finds the lock taken must wait, sometimes by
  entering the kernel (the \code{sys} time grew from 0.001 s to about 0.03 s). The order is the
  same as on the reference machine: \code{race} 0.004 s and wrong, \code{mutex} 0.032 s,
  \code{sem} 0.051 s.
\end{enumerate}

\section{Part 3. A semaphore as a signal}

Here the problem is the order: B must read the message only after A has written it. A semaphore
that starts at 0 is a closed door: B waits at \code{sem\_wait} until A calls \code{sem\_post}.

\begin{lstlisting}[style=ccode,caption={order.c}]
#include <stdio.h>
#include <string.h>
#include <unistd.h>
#include <pthread.h>
#include <semaphore.h>

char message[32];
sem_t ready;

void *writer(void *arg) {
    sleep(1);
    strcpy(message, "hello");
    printf("A: message written\n");
    sem_post(&ready);
    return NULL;
}

void *reader(void *arg) {
    printf("B: waiting for the message...\n");
    sem_wait(&ready);
    printf("B: I read \"%s\"\n", message);
    return NULL;
}

int main(void) {
    sem_init(&ready, 0, 0);
    pthread_t a, b;
    pthread_create(&b, NULL, reader, NULL);
    pthread_create(&a, NULL, writer, NULL);
    pthread_join(a, NULL);
    pthread_join(b, NULL);
    sem_destroy(&ready);
    return 0;
}
\end{lstlisting}

\begin{lstlisting}[caption={order with the semaphore starting at 0}]
$ gcc -Wall -pthread -o order order.c
$ ./order
B: waiting for the message...
A: message written
B: I read "hello"
\end{lstlisting}

B printed its first line at once, nothing happened for one second, then A and B finished.

\begin{lstlisting}[caption={order with sem\_init(\&ready, 0, 1)}]
$ gcc -Wall -pthread -o order order.c && ./order
B: waiting for the message...
B: I read ""
A: message written
\end{lstlisting}

\textbf{Observations:}
\begin{enumerate}
  \item \textbf{Start value 1.} B reads an empty message, \code{B: I read ""}. The semaphore
  already holds one unit, so \code{sem\_wait} passes at once and B reads \code{message} before
  A, which is still sleeping, has written it. The door was open from the start.
  \item \textbf{A semaphore that starts at 0} means ``the event has not happened yet'': whoever
  waits on it is blocked until another thread posts.
\end{enumerate}

\section{Part 4. A semaphore as a counter: the parking lot}

Six cars, three places. The semaphore starts at 3 and counts the free places. The mutex is there
only so that the variable \code{inside} and the \code{printf} stay together.

\begin{lstlisting}[style=ccode,caption={parking.c}]
#include <stdio.h>
#include <unistd.h>
#include <pthread.h>
#include <semaphore.h>

#define PLACES 3
#define CARS 6

sem_t places;
pthread_mutex_t lock = PTHREAD_MUTEX_INITIALIZER;
int inside = 0;

void *car(void *arg) {
    long id = (long)arg;
    sem_wait(&places);
    pthread_mutex_lock(&lock);
    inside++;
    printf("car %ld enters (inside: %d)\n", id, inside);
    pthread_mutex_unlock(&lock);
    sleep(1);
    pthread_mutex_lock(&lock);
    inside--;
    printf("car %ld leaves (inside: %d)\n", id, inside);
    pthread_mutex_unlock(&lock);
    sem_post(&places);
    return NULL;
}

int main(void) {
    sem_init(&places, 0, PLACES);
    pthread_t t[CARS];
    for (long i = 0; i < CARS; i++)
        pthread_create(&t[i], NULL, car, (void *)(i + 1));
    for (int i = 0; i < CARS; i++)
        pthread_join(t[i], NULL);
    sem_destroy(&places);
    return 0;
}
\end{lstlisting}

\begin{lstlisting}[caption={parking with 3 places}]
$ gcc -Wall -pthread -o parking parking.c
$ ./parking
car 2 enters (inside: 1)
car 3 enters (inside: 2)
car 1 enters (inside: 3)
car 2 leaves (inside: 2)
car 3 leaves (inside: 1)
car 1 leaves (inside: 0)
car 4 enters (inside: 1)
car 5 enters (inside: 2)
car 6 enters (inside: 3)
car 4 leaves (inside: 2)
car 6 leaves (inside: 1)
car 5 leaves (inside: 0)
\end{lstlisting}

\begin{lstlisting}[caption={parking with PLACES set to 1}]
$ gcc -Wall -pthread -o parking parking.c && ./parking
car 2 enters (inside: 1)
car 2 leaves (inside: 0)
car 1 enters (inside: 1)
car 1 leaves (inside: 0)
car 3 enters (inside: 1)
car 3 leaves (inside: 0)
car 4 enters (inside: 1)
car 4 leaves (inside: 0)
car 5 enters (inside: 1)
car 5 leaves (inside: 0)
car 6 enters (inside: 1)
car 6 leaves (inside: 0)
\end{lstlisting}

\textbf{Observations:}
\begin{enumerate}
  \item \textbf{The largest value of \code{inside}.} It is 3 and never more. The line that
  guarantees it is \code{sem\_wait(\&places);} together with the start value in
  \code{sem\_init(\&places, 0, PLACES)}: after three cars have taken a unit the value is 0, and
  the fourth car blocks until a car leaves and calls \code{sem\_post}.
  \item \textbf{PLACES set to 1.} The semaphore becomes a mutex (a binary semaphore): only one
  car is inside at a time, each car enters and leaves before the next one, and the run takes
  six seconds instead of two.
\end{enumerate}

\section{Optional: no lock at all}

For a single counter there is a third way: an atomic operation. \code{atomic\_fetch\_add} does
the load, the add and the store as one indivisible step of the CPU.

\begin{lstlisting}[style=ccode,caption={atomic.c}]
#include <stdio.h>
#include <pthread.h>
#include <stdatomic.h>

#define N 1000000

atomic_long counter = 0;

void *add(void *arg) {
    for (int i = 0; i < N; i++)
        atomic_fetch_add(&counter, 1);
    return NULL;
}

int main(void) {
    pthread_t a, b;
    pthread_create(&a, NULL, add, NULL);
    pthread_create(&b, NULL, add, NULL);
    pthread_join(a, NULL);
    pthread_join(b, NULL);
    printf("counter = %ld (expected %d)\n", (long)counter, 2 * N);
    return 0;
}
\end{lstlisting}

\begin{lstlisting}[caption={atomic: correct and faster than the mutex}]
$ gcc -Wall -pthread -o atomic atomic.c
$ time ./atomic
counter = 2000000 (expected 2000000)

real	0m0.016s
user	0m0.029s
sys	0m0.001s
\end{lstlisting}

The answer is correct with no lock, in 0.016 s, half the time of the mutex version, and the
program never enters the kernel to wait.

\plainsection{Conclusions}

The lab showed a race condition with my own eyes: \code{counter++} is three machine steps, and
two threads that interleave those steps lose almost half of their additions, differently at
every run. The fix is to make the three steps a critical section that only one thread can
execute at a time, with a mutex or with a semaphore that starts at 1. The result becomes
correct, at the price of time, because every addition also pays for locking.

The same semaphore then did three different jobs, depending only on its start value. Started at
1 it is a lock. Started at 0 it is a signal that makes one thread wait until another has
finished something, which fixes the order of two actions. Started at \(N\) it is a counter of
free resources, which let at most three cars into the parking lot. A wrong start value breaks
each of them: 2 instead of 1 brings the race back, and 1 instead of 0 lets the reader run too
early.

\end{document}
